\section{总结与延伸}

\subsection{本章小结}
多 GPU 训练从来不是单一技巧，而是通信模式、并行维度和工程约束之间的折中。数据并行适合先跑通系统，张量并行和流水线并行适合把单卡放不下的模型拆开，而 ZeRO/FSDP 则是在训练状态层面继续压缩内存占用。真正的系统设计目标，是让通信代价和模型规模保持同一数量级。

\subsection{拓展阅读}

\begin{itemize}
    \item Megatron-LM 论文：\textit{Efficient Large-Scale Language Model Training on GPU Clusters Using Megatron-LM} (Narayanan et al., 2021)
    \item ZeRO 论文：\textit{ZeRO: Memory Optimizations Toward Training Trillion Parameter Models} (Rajbhandari et al., 2020)
    \item PyTorch FSDP 文档：\url{https://pytorch.org/docs/stable/fsdp.html}
    \item NCCL 性能分析：\url{https://github.com/NVIDIA/nccl-tests/blob/master/doc/PERFORMANCE.md}
    \item GPipe 论文：\textit{GPipe: Efficient Training of Giant Neural Networks using Pipeline Parallelism} (Huang et al., 2019)
    \item PipeDream 论文（1F1B 调度）：\textit{PipeDream: Generalized Pipeline Parallelism for DNN Training} (Narayanan et al., 2019)
\end{itemize}

\end{document}
